Most adversarial attempts to break an aligned language model \emph{fail}, and these ``failures to fail'' are routinely discarded as weak experiments. We argue the opposite: a failed attack can mark a genuine \emph{robustness boundary} that quantifies alignment quality---but only once the behavioral metric is itself validated. We build \ours, a graduated adversarial stress-test harness that applies four failure modes (jailbreak composition, adversarial persona, syntactic/encoding perturbation, and bias elicitation), each with six calibrated intensity levels, identically across five model families spanning three alignment methods. From 2{,}879 real API responses (three seeds, temperature 0.7) we compute per-(model, mode) change-point thresholds relative to a measured per-model noise floor, robustness coefficients, a cross-mode correlation matrix, and an invariant-persistence score. Our central result is a cautionary one: a naive refusal-substring detector scores the entire encoding mode as a universal collapse (compliance $\to 1.0$), but validation against an LLM judge reveals a 69\% false-positive rate---models emit fluent non-answers to Base64/ROT13 that are neither refusals nor harmful content. Judge-corrected, encoding is one of the \emph{most} robust modes (genuine compliance $\approx 0$ at maximum intensity for 4/5 models). The corrected map ranks models \gpt $>$ \gemma $>$ \llama $>$ \qwen $>$ \mistral, shows robustness is multi-dimensional (encoding is orthogonal to an instruction-pressure cluster, $\rho \approx 0$), identifies 11/20 (model, mode) invariants, and finds preference-optimized/deliberative models more robust than fine-tuning-aligned ones (0.77 vs.\ 0.58). Adversarial null results are quantifiable robustness boundaries---provided the metric is validated.
